% ============================================================================
%  explication.tex
%  Document de travail : comprendre la synthèse de la régulation de position
%  (structure LQI, du SISO au MIMO modal), et savoir refaire chaque calcul.
%
%  Compilation : pdflatex explication.tex
%  Ce document est autonome (ne dépend pas de la classe du rapport).
% ============================================================================
\documentclass[a4paper,11pt]{article}

\usepackage[utf8]{inputenc}
\usepackage[T1]{fontenc}
\usepackage[french]{babel}
\usepackage[margin=2.5cm]{geometry}
\usepackage{amsmath, amssymb}
\usepackage{booktabs}
\usepackage{array}
\usepackage{xcolor}
\usepackage[colorlinks=true, linkcolor=blue!50!black]{hyperref}

\newcommand{\qm}{q}
\newcommand{\e}[1]{\cdot 10^{#1}}

\title{Explication --- Régulation de position par retour d'état intégral (LQI)~:\\
du régulateur mono-axe (SISO) à la régulation modale (MIMO)}
\author{Document de travail --- maquette de sustentation magnétique}
\date{\today}

\begin{document}
\maketitle

\tableofcontents
\clearpage

% ============================================================================
\section{Objectif et conventions}
% ============================================================================

Ce document explique, étape par étape, comment sont calculées \textbf{toutes} les
constantes générées par le script \texttt{lqi-inertie.py} et utilisées dans le
firmware (\texttt{T\_MAT}, \texttt{E\_MAT}, \texttt{W\_MAT}, \texttt{F\_STAT},
\texttt{U\_STAT}, \texttt{LQI\_*}, \texttt{OBS\_*}, \texttt{AW\_TOL},
\texttt{REF\_SMOOTH}). L'idée directrice est la suivante~:

\begin{center}
\emph{Il n'existe qu'un seul régulateur, mono-axe. Il est synthétisé une fois
pour le décollage (SISO, un inducteur), puis réutilisé trois fois en vol
(MIMO), une fois par mode de corps rigide~: pompage, tangage, roulis.}
\end{center}

La difficulté n'est donc pas triple~: une fois le régulateur mono-axe compris
(sections~\ref{sec:systeme} à~\ref{sec:marge}), le MIMO n'ajoute que de la
géométrie et de la statique (sections~\ref{sec:modal} à~\ref{sec:synthmimo}).

\paragraph{Conventions.}
\begin{itemize}
    \item L'entrefer $\delta$ \textbf{diminue} quand le corps monte. Une force
          magnétique positive (vers le haut) fait donc diminuer la coordonnée.
          C'est l'origine des signes $-1/m$ et $-1/J$ dans les modèles.
    \item Disposition des inducteurs~: 1 = avant-gauche, 2 = avant-droit,
          3 = arrière-gauche, 4 = arrière-droit. L'axe $x$ est longitudinal
          (positif vers l'avant), l'axe $y$ est latéral (positif à gauche).
    \item Période d'échantillonnage de l'interruption~: $h = 40$~\textmu s
          ($f_e = 25$~kHz).
    \item Les grandeurs sont exprimées en \textbf{écarts} autour du point
          d'équilibre (voir §\ref{subsec:deviation})~; c'est ce qui fait
          disparaître la gravité des modèles dynamiques.
\end{itemize}

% ============================================================================
\section{Le système vu par la boucle de position}
\label{sec:systeme}
% ============================================================================

\subsection{La chaîne d'actionnement~: un seul retard équivalent}
\label{subsec:tauact}

Entre le moment où la régulation de position demande une force $F_c$ et le
moment où cette force agit réellement sur la mécanique, la commande traverse
une cascade de dynamiques rapides~: la boucle de courant fermée, le filtre
coupe-bande, le filtrage analogique et le capteur. Comme dans la synthèse
classique (méthode des petites constantes de temps, cf.\ le bloc $G_{pct}$ du
rapport), on regroupe tout cela en un unique premier ordre~:

\begin{equation}
    G_{act}(s) = \frac{U(s)}{u_c(s)} = \frac{1}{1 + s\,\tau_{act}},
    \qquad
    \tau_{act} = \tau_{iBF} + \tau_{notch} + \tau_{RC} + \tau_{mes}.
\end{equation}

\begin{table}[h]
    \centering
    \begin{tabular}{llr}
        \toprule
        Contribution & Origine & Valeur \\
        \midrule
        $\tau_{iBF}$   & boucle de courant fermée (moyenne des 4 inducteurs) & $1{,}298$~ms \\
        $\tau_{notch}$ & retard équivalent du filtre coupe-bande 75~Hz       & $3{,}98$~ms  \\
        $\tau_{RC}$    & filtre analogique $RC = 4{,}3\,\text{k}\Omega \times 33\,\text{nF}$ & $141{,}9$~\textmu s \\
        $\tau_{mes}$   & capteur de position, $f_{m\delta} = 1{,}1$~kHz~: $1/(2\pi f_{m\delta})$ & $144{,}7$~\textmu s \\
        \midrule
        $\tau_{act}$   & somme & $\mathbf{5{,}565}$~\textbf{ms} \\
        \bottomrule
    \end{tabular}
    \caption{Constantes de temps regroupées dans $\tau_{act}$.}
    \label{tab:tauact}
\end{table}

Deux remarques importantes pour le rapport~:
\begin{itemize}
    \item $\tau_{mes} = 144{,}7$~\textmu s et $\tau_{notch} = 3{,}98$~ms sont
          exactement les valeurs déjà utilisées dans le chapitre de régulation
          ($\tau_{m\delta}$ et $\tau_{F_C}$).
    \item $\tau_{iBF} \approx 1{,}3$~ms est nettement plus lent que la valeur
          théorique de $78$~\textmu s calculée avec $K_{pi} = 91$. C'est la
          conséquence directe du choix expérimental $K_{pi} \approx 2$~: la
          boucle de courant réellement implantée est plus lente, et c'est
          \emph{cette} valeur qu'il faut utiliser pour la synthèse de position.
\end{itemize}

C'est la valeur $1/\tau_{act} \approx 180$~rad/s qui borne la rapidité que
l'on peut demander à la boucle de position~: on ne peut pas placer des pôles
de régulation plus rapides que la chaîne d'actionnement elle-même.

\subsection{La mécanique en variables d'écart}
\label{subsec:deviation}

Pour un axe de masse (ou d'inertie) $J$, l'équation de Newton s'écrit, avec la
convention de signe de l'entrefer~:
\begin{equation}
    J\,\ddot{\delta} = J\,g_{eq} - F,
\end{equation}
où $F$ est la force (ou le couple) appliqué et $J g_{eq}$ le terme de pesanteur.
À l'équilibre, $F_0 = J g_{eq}$ (le poids est exactement compensé). En posant
les \textbf{variables d'écart} $\qm = \delta - \delta_0$ et $U = F - F_0$~:
\begin{equation}
    \boxed{\;\ddot{\qm} = -\frac{1}{J}\,U\;}
    \label{eq:doubleint}
\end{equation}
La gravité a disparu~: elle est prise en charge séparément par un terme
\emph{statique} ajouté à la commande (le \texttt{FP} du SISO, le
\texttt{F\_STAT} du MIMO, voir §\ref{subsec:statique}). Le système à régler
est un \textbf{double intégrateur}~: c'est le même objet pour le pompage
($J = m$), le tangage et le roulis ($J$ = moment d'inertie), et pour un
inducteur seul ($J = m_{tot}/4 = 4{,}513$~kg).

% ============================================================================
\section{Modèle d'état d'un axe}
\label{sec:etat}
% ============================================================================

On empile quatre états~: la position d'écart $\qm$, la vitesse $\dot{\qm}$,
l'effort réellement appliqué $U$ (sortie du premier ordre $G_{act}$), et
l'état intégral $\varepsilon$~:
\begin{equation}
    x = \begin{bmatrix} \qm \\ \dot{\qm} \\ U \\ \varepsilon \end{bmatrix},
    \qquad
    \left\{
    \begin{aligned}
        \dot{\qm}          &= \dot{\qm} \\
        \ddot{\qm}         &= -\tfrac{1}{J}\,U \\
        \dot{U}            &= \tfrac{1}{\tau_{act}}\,(u_c - U) \\
        \dot{\varepsilon}  &= \qm_c - \qm \;\;(= -\qm \text{ en écart, consigne nulle})
    \end{aligned}
    \right.
\end{equation}
soit sous forme matricielle $\dot{x} = A\,x + B\,u_c$~:
\begin{equation}
    A = \begin{bmatrix}
        0 & 1 & 0 & 0 \\
        0 & 0 & -\frac{1}{J} & 0 \\
        0 & 0 & -\frac{1}{\tau_{act}} & 0 \\
        -1 & 0 & 0 & 0
    \end{bmatrix},
    \qquad
    B = \begin{bmatrix} 0 \\ 0 \\ \frac{1}{\tau_{act}} \\ 0 \end{bmatrix}.
    \label{eq:AB}
\end{equation}

C'est \emph{exactement} le vecteur d'état $x = [x_{em},\ \delta,\ \dot\delta,\
x_r]$ du chapitre de régulation (à l'ordre des composantes près)~: $U$ joue le
rôle de $x_{em}$ (l'effort filtré par $G_{pct}$) et $\varepsilon$ celui de
$x_r$ (l'intégrateur). Rien de nouveau~: seule la méthode de calcul des gains
change.

% ============================================================================
\section{La loi de commande~: structure LQI}
\label{sec:lqi}
% ============================================================================

\subsection{Retour d'état avec action intégrale}

La commande est un retour d'état sur le système augmenté~:
\begin{equation}
    u_c = -K\,x = -\big(k_q\,\qm + k_{\dot q}\,\dot{\qm} + k_U\,U
          + k_\varepsilon\,\varepsilon\big).
    \label{eq:loi}
\end{equation}
Chaque terme a un rôle physique clair~:
\begin{itemize}
    \item $k_q\,\qm$~: le ressort. Plus l'écart de position est grand, plus on
          pousse. Fixe la raideur de la lévitation.
    \item $k_{\dot q}\,\dot{\qm}$~: l'amortisseur. Freine le mouvement, apporte
          l'amortissement qu'un double intégrateur n'a pas du tout. Sans lui,
          le système oscille indéfiniment.
    \item $k_\varepsilon\,\varepsilon$~: l'intégrale de l'erreur. Élimine
          l'erreur statique due aux imperfections du modèle (constante de
          force réelle, offsets, masse mal connue). C'est le \og I \fg{} de LQI.
    \item $k_U\,U$~: retour sur l'effort filtré. Ce terme est
          \textbf{volontairement omis dans le firmware} (voir
          §\ref{subsec:ku}).
\end{itemize}

Dans le firmware, la loi \eqref{eq:loi} apparaît sous la forme
\begin{equation}
    u_{cmd} = Q\,(\qm - \qm_{ref}) + Q_D\,\hat{\dot q} - E_{PS}\,\varepsilon,
    \qquad \varepsilon \mathrel{+}= (\qm_{ref} - \qm)\cdot h,
\end{equation}
avec $Q = |k_q|$, $Q_D = |k_{\dot q}|$, $E_{PS} = |k_\varepsilon|$
(\texttt{LQI\_Q}, \texttt{LQI\_QD}, \texttt{LQI\_EPS}). Les valeurs absolues
sont légitimes car les signes sortent toujours identiques du placement avec
notre convention d'entrefer ($k_q, k_{\dot q} < 0$ et $k_\varepsilon > 0$,
les signes étant réabsorbés dans l'écriture ci-dessus). La vitesse
$\hat{\dot q}$ n'est pas mesurée~: elle provient de l'observateur
(section~\ref{sec:obs}).

\subsection{Pourquoi \og LQI \fg{}, et le lien avec Riccati}
\label{subsec:pourquoilqi}

Le nom \textbf{LQI} (\emph{Linear Quadratic Integral}) désigne la
\emph{structure}~: retour d'état sur le système augmenté d'un intégrateur
d'erreur. Historiquement, les gains $K$ d'une telle structure sont obtenus en
minimisant le critère quadratique
\begin{equation}
    J_{LQ} = \int_0^{\infty} \big( x^T Q_{LQ}\, x + u^T R_{LQ}\, u \big)\,dt,
\end{equation}
où $Q_{LQ}$ pénalise les écarts d'état et $R_{LQ}$ l'énergie de commande~; la
solution passe par l'équation algébrique de Riccati (fonctions
\texttt{lqr}/\texttt{lqi} de Matlab, \texttt{scipy}).

\textbf{Point important pour le rapport}~: dans notre cas, les gains ne sont
\emph{pas} obtenus par Riccati mais par \textbf{placement de pôles} sur le
même système augmenté (section~\ref{sec:placement}). Les deux méthodes
produisent la même structure de commande~\eqref{eq:loi}~; seule la façon de
choisir $K$ diffère~:
\begin{itemize}
    \item placement de pôles~: on impose directement la dynamique de boucle
          fermée (les pôles), méthode identique à celle du chapitre de
          régulation (identification du polynôme caractéristique)~;
    \item LQ/Riccati~: on impose des pondérations physiques et l'algorithme en
          déduit les pôles.
\end{itemize}
Il est donc plus rigoureux d'écrire dans le rapport~: \og régulateur par
retour d'état avec action intégrale (structure LQI), gains obtenus par
placement de pôles \fg. Si un expert demande \og quelles matrices $Q_{LQ}$ et
$R_{LQ}$ avez-vous choisies~? \fg, la réponse honnête est~: aucune, le réglage
est fait par placement de pôles, validé ensuite par simulation et par un
critère de marge (section~\ref{sec:marge}).

% ============================================================================
\section{Discrétisation exacte à 25 kHz}
\label{sec:disc}
% ============================================================================

Le régulateur tourne dans une interruption à $f_e = 25$~kHz. La synthèse est
donc faite directement en \textbf{temps discret}. Pour un système continu
$\dot x = Ax + Bu$ maintenu par bloqueur d'ordre zéro (la commande est
constante entre deux interruptions), la discrétisation \emph{exacte} est~:
\begin{equation}
    x_{k+1} = A_d\,x_k + B_d\,u_k,
    \qquad
    A_d = e^{A h},
    \qquad
    B_d = \int_0^{h} e^{A\sigma}\,d\sigma \; B .
\end{equation}
En pratique on calcule les deux d'un coup par l'astuce de la matrice
augmentée~:
\begin{equation}
    \exp\!\left( \begin{bmatrix} A & B \\ 0 & 0 \end{bmatrix} h \right)
    = \begin{bmatrix} A_d & B_d \\ 0 & I \end{bmatrix},
\end{equation}
ce qui est exactement la fonction \texttt{c2d()} du script.

\paragraph{Formes fermées (calcul à la main).}
Pour notre $A$ \eqref{eq:AB}, l'exponentielle se calcule analytiquement, et
comme $h \ll \tau_{act}$ les termes utiles se réduisent à~:
\begin{equation}
    \begin{aligned}
        &\text{position} & \qm_{k+1} &\approx \qm_k + h\,\dot{\qm}_k
            &&(\text{terme } h = 4\e{-5}) \\
        &\text{vitesse} & \dot{\qm}_{k+1} &\approx \dot{\qm}_k
            - \frac{\tau_{act}}{J}\big(1 - e^{-h/\tau_{act}}\big)\,U_k
            \;\approx\; \dot{\qm}_k - \frac{h}{J}\,U_k \\
        &\text{effort} & U_{k+1} &= e^{-h/\tau_{act}}\,U_k
            + \big(1 - e^{-h/\tau_{act}}\big)\,u_{c,k} \\
        &\text{intégrale} & \varepsilon_{k+1} &\approx \varepsilon_k
            + h\,(\qm_c - \qm_k)
    \end{aligned}
    \label{eq:discret}
\end{equation}
Numériquement, avec $h = 40$~\textmu s et $\tau_{act} = 5{,}565$~ms~:
\begin{equation}
    e^{-h/\tau_{act}} = 0{,}992\,837, \qquad 1 - e^{-h/\tau_{act}} = 0{,}007\,163 .
\end{equation}
On reconnaît immédiatement les constantes \texttt{OBS\_AD33} et
\texttt{OBS\_BD3} du firmware~: ce sont ces exponentielles-là.

\paragraph{Correspondance des pôles.}
Un pôle continu $s = p$ devient en discret $z = e^{p h}$. Les pôles de
synthèse choisis (voir §\ref{subsec:choixpoles}) donnent~:
\begin{equation}
    \{-80, -85, -90, -20\}\ \text{rad/s}
    \;\longrightarrow\;
    \{0{,}996\,805,\; 0{,}996\,606,\; 0{,}996\,406,\; 0{,}999\,200\}.
\end{equation}
Des pôles discrets très proches de 1~: c'est normal, la boucle de position est
très lente devant $f_e$.

% ============================================================================
\section{Calcul des gains par placement de pôles}
\label{sec:placement}
% ============================================================================

\subsection{Principe}

En boucle fermée, $u_k = -K x_k$ donne $x_{k+1} = (A_d - B_d K)\,x_k$. On
cherche $K = [k_q,\ k_{\dot q},\ k_U,\ k_\varepsilon]$ tel que~:
\begin{equation}
    \det\big(z I - (A_d - B_d K)\big)
    \;=\; \prod_{i=1}^{4} \big(z - e^{p_i h}\big).
    \label{eq:place}
\end{equation}
C'est un système de 4 équations (les coefficients du polynôme en $z^3, z^2,
z^1, z^0$) à 4 inconnues (les gains) --- \textbf{exactement la démarche du
chapitre de régulation}, où le déterminant est développé à la main en continu
puis identifié au polynôme cible. La seule différence est qu'ici le calcul est
fait en discret et résolu numériquement (\texttt{place\_poles} de
\texttt{scipy}, équivalent de \texttt{place} de Matlab). Le développement à la
main reste possible (même technique de déterminant $4\times4$), simplement
fastidieux en discret~; pour un calcul manuel, on peut faire le placement en
continu comme dans le rapport puis vérifier que $h$ est assez petit pour que
le résultat coïncide (à 25~kHz, l'écart est négligeable).

\subsection{Choix des pôles}
\label{subsec:choixpoles}

\begin{itemize}
    \item \textbf{Trois pôles \og rapides \fg} $\{-80, -85, -90\}$~rad/s~:
          la dynamique de régulation proprement dite ($\approx 13$--$14$~Hz).
          Ils restent bien en dessous de la limite de la chaîne d'actionnement
          $1/\tau_{act} \approx 180$~rad/s.
    \item \textbf{Un pôle lent} $-20$~rad/s pour la branche intégrale~:
          l'action intégrale doit être plus lente que la régulation qu'elle
          corrige, sinon elle dégrade l'amortissement (c'est le rôle du
          rapport $T_i$ dans la structure classique).
    \item Les pôles sont \textbf{légèrement distincts} (80, 85, 90) plutôt que
          confondus~: exigence de l'algorithme de placement (pôles multiples
          mal conditionnés), sans effet pratique sur la réponse.
\end{itemize}
Le choix précis de ces valeurs est itératif~: placement $\to$ simulation
$\to$ critère de marge (section~\ref{sec:marge}) $\to$ essai en vol. Les
commentaires du script gardent la trace de cette itération (pôles ralentis
après l'observation d'oscillations croissantes en vol, dues à une marge trop
faible en saturation).

\subsection{Résultats numériques}

Pour le régulateur SISO du décollage ($J = m = 4{,}513$~kg, pôles
$\{-80,-85,-90,-95\}$) et pour les trois modes MIMO (pôles
$\{-80,-85,-90,-20\}$, inerties de la section~\ref{sec:inerties})~:

\begin{table}[h]
    \centering
    \begin{tabular}{lrrrrr}
        \toprule
        Axe & $J$ & $Q = |k_q|$ & $Q_D = |k_{\dot q}|$ & $E_{PS} = |k_\varepsilon|$ & $k_U$ (omis) \\
        \midrule
        SISO décollage & $4{,}513$~kg & $66\,822$ & $1\,149{,}5$ & $1\,455\,109$ & $0{,}95$ \\
        Z (pompage)    & $15{,}697$~kg & $91\,128$ & $2\,334{,}0$ & $1\,067\,127$ & $0{,}53$ \\
        T (tangage)    & $0{,}7373$~kg$\cdot$m$^2$ & $4\,280{,}5$ & $109{,}6$ & $50\,125$ & $0{,}53$ \\
        R (roulis)     & $0{,}7373$~kg$\cdot$m$^2$ & $4\,280{,}5$ & $109{,}6$ & $50\,125$ & $0{,}53$ \\
        \bottomrule
    \end{tabular}
    \caption{Gains obtenus par placement de pôles.}
    \label{tab:gains}
\end{table}

Une propriété remarquable et facile à vérifier à la main~: \textbf{à pôles
identiques, les gains sont proportionnels à l'inertie}. En effet le seul
paramètre du modèle qui contient $J$ est le gain $1/J$ de \eqref{eq:doubleint},
donc doubler $J$ exige de doubler tous les efforts pour la même trajectoire.
Vérification~: $Q_Z / Q_T = 91\,128 / 4\,280{,}5 = 21{,}29 = J_Z / J_T =
15{,}697 / 0{,}7373$. Les trois modes ont ainsi \emph{exactement la même
dynamique de boucle fermée}, seuls les niveaux d'effort changent.

\subsection{L'omission du gain $k_U$}
\label{subsec:ku}

Le placement fournit quatre gains, mais le firmware n'en applique que trois~:
le retour sur l'effort filtré $U$ (valeur $k_U \approx 0{,}5$--$0{,}9$) est
ignoré. Deux justifications~:
\begin{enumerate}
    \item $U$ n'est pas mesurable directement (il faudrait le reconstruire via
          l'observateur, ce qui ajouterait du bruit de commande pour un gain
          faible)~;
    \item c'est la même simplification que dans la structure du rapport, où le
          vecteur de gains s'écrit $K = [0,\ K_\delta,\ K_{\dot\delta},\ -K_r]$
          avec un \emph{zéro} sur $x_{em}$.
\end{enumerate}
Conséquence honnête à assumer~: les pôles réellement obtenus ne sont pas
exactement ceux placés. C'est pourquoi le script re-simule ensuite la boucle
\emph{telle qu'implantée} (trois gains, observateur, filtre coupe-bande,
saturations) et vérifie la stabilité et la fréquence dominante --- c'est le
rôle de \texttt{simulate\_mode()}.

% ============================================================================
\section{L'observateur de Luenberger}
\label{sec:obs}
% ============================================================================

\subsection{Pourquoi un observateur}

La loi \eqref{eq:loi} a besoin de la vitesse $\dot{\qm}$, qui n'est pas
mesurée. Dériver numériquement la position à 25~kHz est inutilisable~: le
bruit de quantification de l'ADC serait amplifié d'un facteur $1/h = 25\,000$.
L'observateur remplace la dérivation par une \emph{prédiction par le modèle,
recalée par la mesure}.

\subsection{Structure et équations}

On observe les trois états physiques (l'intégrateur $\varepsilon$ est connu
exactement, inutile de l'observer)~:
\begin{equation}
    x_{obs} = \begin{bmatrix} \qm \\ \dot{\qm} \\ U \end{bmatrix},
    \qquad
    A_{obs} = \begin{bmatrix}
        0 & 1 & 0 \\ 0 & 0 & -\frac{1}{J} \\ 0 & 0 & -\frac{1}{\tau_{act}}
    \end{bmatrix},
    \quad
    B_{obs} = \begin{bmatrix} 0 \\ 0 \\ \frac{1}{\tau_{act}} \end{bmatrix},
    \quad
    C_{obs} = \begin{bmatrix} 1 & 0 & 0 \end{bmatrix}.
\end{equation}
Après discrétisation (mêmes formules qu'en section~\ref{sec:disc}),
l'observateur s'écrit~:
\begin{equation}
    \hat{x}_{k+1} = A_{d,obs}\,\hat{x}_k + B_{d,obs}\,\bar u_k
                    + L\,\underbrace{(\qm_{mes,k} - \hat{\qm}_k)}_{e_k},
    \label{eq:obs}
\end{equation}
où $\bar u$ est l'effort \emph{réellement appliqué} (après saturation et
filtre), et $L = [L_1, L_2, L_3]^T$ les gains de recalage. Développé ligne à
ligne, c'est mot pour mot le code de l'interruption~:
\begin{equation}
    \begin{aligned}
        \hat{\qm}^{+}      &= \hat{\qm} + \underbrace{h}_{\texttt{AD12}}\,\hat{\dot q} + L_1\,e \\
        \hat{\dot q}^{+}   &= \hat{\dot q} + \underbrace{\Big(-\tfrac{\tau_{act}}{J}\big(1 - e^{-h/\tau_{act}}\big)\Big)}_{\texttt{AD23}\;\approx\; -h/J}\,\hat{U} + L_2\,e \\
        \hat{U}^{+}        &= \underbrace{e^{-h/\tau_{act}}}_{\texttt{AD33}}\,\hat{U} + \underbrace{\big(1 - e^{-h/\tau_{act}}\big)}_{\texttt{BD3}}\,\bar u + L_3\,e
    \end{aligned}
\end{equation}
Les termes croisés d'ordre supérieur de l'exponentielle (par exemple l'entrée
$(1,3)$ de $A_{d,obs}$, en $h^2/2J$) sont négligés~: à $h = 40$~\textmu s ils
sont de l'ordre de $10^{-11}$.

\subsection{Calcul des gains $L$ par dualité}

L'erreur d'estimation $e_k = x_k - \hat{x}_k$ obéit à
$e_{k+1} = (A_{d,obs} - L\,C_{obs})\,e_k$~: on veut qu'elle converge vite,
donc on \emph{place les pôles} de $A_{d,obs} - L C_{obs}$. Or les valeurs
propres d'une matrice et de sa transposée sont identiques~:
\begin{equation}
    \operatorname{eig}(A - LC) = \operatorname{eig}(A^T - C^T L^T),
\end{equation}
si bien que trouver $L$ est \emph{le même problème} que trouver un retour
d'état $K$ pour la paire $(A^T, C^T)$ --- c'est la \textbf{dualité
commande/observation}, et la raison de la ligne
\texttt{place\_poles(Ad\_obs.T, C\_obs.T, ...)} du script.

Pôles d'observation choisis~: $\{-700, -760, -820\}$~rad/s, environ 8 à 10
fois plus rapides que la régulation. Règle de réglage~:
\begin{itemize}
    \item trop lent~: la boucle de commande s'appuie sur des estimations en
          retard $\to$ perte d'amortissement~;
    \item trop rapide~: $L$ devient grand et recopie le bruit de mesure dans
          $\hat{\dot q}$, donc dans la commande.
\end{itemize}
Valeurs obtenues (identiques pour tous les axes en $L_1$, $L_2$ car la
géométrie de la mesure est la même~; $L_3$ dépend de $J$)~:
$L_1 = 0{,}0827$, $L_2 = 52{,}15$, et $L_3 = -115\,565$ (Z), $-5\,428$ (T/R),
$-33\,225$ (SISO).

% ============================================================================
\section{Marge de stabilité et saturations}
\label{sec:marge}
% ============================================================================

\subsection{Le critère $\omega^{*}$ (analyse à la main)}

Vue de l'extérieur, la loi de commande est un PID~: en négligeant
l'observateur, la fonction de transfert du régulateur est
\begin{equation}
    C(j\omega) = Q + j\left(Q_D\,\omega - \frac{E_{PS}}{\omega}\right).
\end{equation}
Il existe une pulsation où les effets dérivé et intégral se compensent
exactement et où le régulateur est \emph{purement proportionnel}~:
\begin{equation}
    \boxed{\;\omega^{*} = \sqrt{\frac{E_{PS}}{Q_D}}\;}
\end{equation}
À cette pulsation, la phase de la boucle ouverte est celle du double
intégrateur, soit $-180^\circ$ (le retard $\tau_{act}$ n'ajoute que
$\arctan(\omega^*\tau_{act}) \approx 7^\circ$, négligé ici). Le module de la
boucle ouverte y vaut~:
\begin{equation}
    \boxed{\;\big|L(j\omega^{*})\big| = \frac{Q}{J\,\omega^{*2}}
    = \frac{Q\,Q_D}{J\,E_{PS}}\;}
    \label{eq:marge}
\end{equation}

\textbf{Pourquoi ce nombre doit être grand.} En dessous de $\omega^{*}$,
l'intégrale domine et la phase descend vers $-270^\circ$~: la boucle est
\emph{conditionnellement stable}. Le diagramme de Nyquist passe du bon côté du
point critique uniquement si le gain à $\omega^{*}$ est \emph{supérieur} à 1.
Autrement dit, contrairement à une marge de gain classique, c'est une
\emph{baisse} de gain qui déstabilise. Le rapport \eqref{eq:marge} est donc le
facteur de réduction de gain admissible.

\textbf{Le lien avec la saturation.} Quand les forces saturent (butées
$0 \ldots F_{max}$), tout se passe comme si le gain effectif de la boucle
diminuait (point de vue \og premier harmonique \fg~: la sortie écrêtée
contient moins de fondamental que la commande demandée). Si
$|L(j\omega^{*})|$ est proche de 1, une grosse perturbation qui provoque des
saturations amène la boucle exactement sur son point critique~: on observe
alors une oscillation quasi sinusoïdale à $f^{*} = \omega^{*}/2\pi$ qui
croît lentement --- c'est précisément le phénomène observé en vol et documenté
dans les commentaires du script (points critiques mesurés à $3{,}3$ et
$14{,}8$~Hz), corrigé en ralentissant les pôles pour remonter la marge.

Valeurs de la synthèse actuelle~:
\begin{equation}
    \omega^{*} = 21{,}4\ \text{rad/s} \;(3{,}40\ \text{Hz}),
    \qquad
    |L(j\omega^{*})| = 12{,}7 \;(\approx 22\ \text{dB})
    \quad\text{pour les trois modes.}
\end{equation}
Le script exige $|L(j\omega^{*})| > 8$ et affiche un avertissement sinon.

\subsection{Anti-windup par intégration conditionnelle}

Pendant une saturation, l'erreur persiste et l'intégrateur, s'il continuait à
intégrer, accumulerait une valeur énorme (\emph{windup}) qu'il faudrait
ensuite \og désapprendre \fg, au prix de gros dépassements. Le firmware
applique une \textbf{intégration conditionnelle}~: à chaque période, on
compare l'effort demandé $u_{cmd}$ à l'effort réellement réalisable $\bar u$
(reconstruit à partir des forces saturées, voir §\ref{subsec:antiwindupmimo}),
et l'on ne met à jour $\varepsilon$ que si l'écart est inférieur à la
tolérance \texttt{AW\_TOL} (2\,\% de l'effort statique).

% ============================================================================
\section{Du SISO au MIMO~: les coordonnées modales}
\label{sec:modal}
% ============================================================================

\subsection{Pourquoi changer de coordonnées}

Le véhicule est un corps rigide~: si l'inducteur 1 tire plus fort, il ne
modifie pas seulement $\delta_1$ mais, par le couple créé, aussi $\delta_2,
\delta_3, \delta_4$. Quatre boucles SISO indépendantes se \og disputent \fg{}
donc le même corps. Un corps rigide ne possède ici que \textbf{trois degrés de
liberté}~: la translation verticale $z$ (\textbf{pompage} --- on ne dit pas
\og hauteur \fg{} car on désigne le \emph{mode de mouvement} du centre de
gravité, pas une mesure en un point), le \textbf{tangage} $\theta$ (rotation
autour de l'axe transversal) et le \textbf{roulis} $\varphi$ (rotation autour
de l'axe longitudinal). Dans ces coordonnées, la dynamique se
\textbf{découple}~: la force totale n'agit que sur $z$, le couple de tangage
que sur $\theta$, le couple de roulis que sur $\varphi$. Le MIMO devient trois
SISO.

\subsection{Cinématique~: la matrice $T$ (calcul à la main)}
\label{subsec:T}

Positions des inducteurs par rapport au centre de gravité (empattement
longitudinal $28$~cm, voie latérale $26$~cm)~:
\begin{equation}
    x_I = [+0{,}14,\ +0{,}14,\ -0{,}14,\ -0{,}14]\ \text{m},
    \qquad
    y_I = [+0{,}13,\ -0{,}13,\ +0{,}13,\ -0{,}13]\ \text{m}.
\end{equation}
Pour de petites rotations, l'entrefer au droit de l'inducteur $k$ vaut~:
\begin{equation}
    \delta_k \approx z + x_k\,\theta + y_k\,\varphi
    \qquad\Longleftrightarrow\qquad
    \boldsymbol{\delta} = G\,q,
    \quad
    G = \begin{bmatrix}
        1 & +0{,}14 & +0{,}13 \\
        1 & +0{,}14 & -0{,}13 \\
        1 & -0{,}14 & +0{,}13 \\
        1 & -0{,}14 & -0{,}13
    \end{bmatrix},
    \quad
    q = \begin{bmatrix} z \\ \theta \\ \varphi \end{bmatrix}.
\end{equation}
$G$ est $4\times3$~: on a 4 mesures pour 3 inconnues (une mesure redondante,
la combinaison $\delta_1 - \delta_2 - \delta_3 + \delta_4$, dite de
\emph{torsion}, est nulle pour un corps rigide). On récupère $q$ au sens des
\textbf{moindres carrés}~:
\begin{equation}
    q = \underbrace{(G^T G)^{-1} G^T}_{T}\,\boldsymbol{\delta}.
\end{equation}
Le calcul est trivial à la main car la géométrie est symétrique
($\sum x_k = \sum y_k = \sum x_k y_k = 0$), donc $G^T G$ est
\textbf{diagonale}~:
\begin{equation}
    G^T G = \operatorname{diag}\Big(4,\ \textstyle\sum x_k^2,\ \sum y_k^2\Big)
          = \operatorname{diag}\big(4,\ 0{,}0784,\ 0{,}0676\big),
\end{equation}
d'où les trois lignes de $T$, qui se lisent physiquement~:
\begin{equation}
    z = \frac{\delta_1+\delta_2+\delta_3+\delta_4}{4},
    \qquad
    \theta = \frac{(\delta_1+\delta_2)-(\delta_3+\delta_4)}{4 \times 0{,}14},
    \qquad
    \varphi = \frac{(\delta_1+\delta_3)-(\delta_2+\delta_4)}{4 \times 0{,}13}.
\end{equation}
Soit numériquement $T_{2k} = \pm 1{,}7857$ et $T_{3k} = \pm 1{,}9231$~:
exactement les valeurs de \texttt{T\_MAT}. Le pompage est la \emph{moyenne}
des entrefers, le tangage l'écart avant/arrière divisé par le bras de levier,
le roulis l'écart gauche/droite divisé par le sien.

\subsection{Statique~: les matrices $E$ et $W$}
\label{subsec:statique}

\paragraph{Des forces aux efforts modaux ($E$).}
La résultante et les moments des 4 forces sont~:
\begin{equation}
    u = \begin{bmatrix} \sum_k F_k \\ \sum_k x_k F_k \\ \sum_k y_k F_k \end{bmatrix}
      = \underbrace{G^T}_{E}\,F .
\end{equation}
$E = G^T$~: ligne 1 = force totale (pompage), ligne 2 = couple de tangage,
ligne 3 = couple de roulis. Aucun calcul~: c'est la définition des moments.

\paragraph{Des efforts modaux aux forces (allocation $W$).}
Problème inverse~: le régulateur demande 3 efforts $u$, il faut trouver 4
forces telles que $E F = u$. Système sous-déterminé (4 inconnues, 3
équations)~: on choisit la solution de \textbf{norme minimale}, donnée par la
pseudo-inverse à droite~:
\begin{equation}
    W = E^T (E E^T)^{-1} = G\,(G^T G)^{-1}.
\end{equation}
Grâce à la diagonalité de $G^T G$, chaque force se calcule à la main~:
\begin{equation}
    \boxed{\;F_k = \frac{u_Z}{4} + \frac{x_k}{\sum x^2}\,u_T
    + \frac{y_k}{\sum y^2}\,u_R\;}
    \qquad\text{soit}\qquad
    F_k = \frac{u_Z}{4} \pm 1{,}7857\,u_T \pm 1{,}9231\,u_R .
\end{equation}
Deux propriétés à connaître (et à citer devant un expert)~:
\begin{itemize}
    \item $E\,W = I_3$~: ce qu'on demande est ce qu'on obtient (hors
          saturation). C'est le \texttt{assert} du script.
    \item $W = T^T$~: l'allocation est la transposée de la reconstruction.
          Ce n'est pas un hasard, c'est la dualité cinématique/statique
          (principe des travaux virtuels)~: les mêmes bras de levier servent
          dans les deux sens.
    \item Le \textbf{noyau de $E$} (la combinaison de forces
          $[+1,-1,-1,+1]/2$, la \og torsion \fg) ne produit ni force ni couple~:
          l'allocation par pseudo-inverse ne l'excite jamais. C'est la
          résolution propre de la redondance 4 actionneurs / 3 degrés de
          liberté.
\end{itemize}

\paragraph{Compensation du poids (\texttt{U\_STAT}, \texttt{F\_STAT}).}
L'effort statique à fournir est le poids, sans couples (centre de gravité
quasi centré, vérifié à moins de 3~mm par la méthode du
§\ref{subsec:identcg})~:
\begin{equation}
    U_{stat} = \begin{bmatrix} m_{tot}\,g / K_{FF} \\ 0 \\ 0 \end{bmatrix}
             = \begin{bmatrix} 154{,}0 \\ 0 \\ 0 \end{bmatrix}\ \text{N},
    \qquad
    F_{stat} = W\,U_{stat}
             = \frac{m_{tot}\,g}{4\,K_{FF}}\begin{bmatrix}1\\1\\1\\1\end{bmatrix}
             = 38{,}5\ \text{N par inducteur}.
\end{equation}
Le facteur $K_{FF} = 1{,}15$ est un \textbf{calage expérimental}~: la
constante de force théorique $k$ surestime la force réelle (fuites, franges,
saturation du fer --- l'identification en vol donne $k_{r\acute{e}el} \approx
0{,}64\,k_{th}$). Plutôt que de corriger $k$ partout, le feedforward est
volontairement réduit et \textbf{l'intégrateur du mode Z rattrape l'écart}~:
c'est exactement son rôle. C'est un point de conception à assumer tel quel
dans le rapport.

\subsection{Dynamique modale}

Newton--Euler sur le corps rigide, en variables d'écart et avec la convention
d'entrefer (§\ref{subsec:deviation})~:
\begin{equation}
    m_{tot}\,\ddot z = -u_Z,
    \qquad
    J_\theta\,\ddot\theta = -u_T,
    \qquad
    J_\varphi\,\ddot\varphi = -u_R .
\end{equation}
Chaque mode est un double intégrateur \eqref{eq:doubleint}~: on applique
\emph{trois fois} la synthèse des sections~\ref{sec:etat}
à~\ref{sec:obs}, en remplaçant simplement $J$ par $m_{tot}$, $J_\theta$ ou
$J_\varphi$.

% ============================================================================
\section{Les inerties~: mesure au pendule bifilaire et calage en vol}
\label{sec:inerties}
% ============================================================================

\subsection{Le pendule bifilaire}

Pour mesurer un moment d'inertie sans démonter la maquette, on la suspend par
\textbf{deux fils verticaux} parallèles de longueur $L$, écartés de $2B$
symétriquement autour du centre de gravité, et on la fait osciller en rotation
autour de l'axe vertical passant par le centre.

\emph{Mise en équation (petites oscillations)~:} pour une rotation $\psi$,
chaque point d'attache se déplace horizontalement de $B\psi$~; chaque fil
s'incline donc de $B\psi/L$ et sa tension ($m g/2$ par fil) acquiert une
composante horizontale $\tfrac{mg}{2}\cdot\tfrac{B\psi}{L}$ qui s'oppose à la
rotation avec un bras de levier $B$. Le couple de rappel total vaut~:
\begin{equation}
    C_r = -2 \cdot \frac{m g}{2}\cdot\frac{B\psi}{L}\cdot B
        = -\frac{m g B^2}{L}\,\psi
    \qquad\Longrightarrow\qquad
    J\ddot\psi = -\frac{m g B^2}{L}\,\psi .
\end{equation}
C'est un oscillateur harmonique de période $T_p = 2\pi\sqrt{J L /(m g B^2)}$,
d'où, en mesurant $T_p$ au chronomètre~:
\begin{equation}
    \boxed{\;J = \frac{m\,g\,B^2\,T_p^2}{4\pi^2 L}\;}
\end{equation}

Application avec $L = 1{,}23$~m et $m_{tot} = 18{,}052$~kg~:
\begin{align}
    J_X &= \frac{18{,}052 \cdot 9{,}81 \cdot 0{,}1925^2 \cdot 2{,}267^2}{4\pi^2 \cdot 1{,}23}
         = 0{,}695\ \text{kg}\cdot\text{m}^2
         &&(2B = 0{,}385\ \text{m},\ T_p = 2{,}267\ \text{s}),\\
    J_Y &= \frac{18{,}052 \cdot 9{,}81 \cdot 0{,}139^2 \cdot 3{,}833^2}{4\pi^2 \cdot 1{,}23}
         = 1{,}035\ \text{kg}\cdot\text{m}^2
         &&(2B = 0{,}278\ \text{m},\ T_p = 3{,}833\ \text{s}).
\end{align}
\emph{Attention à décrire honnêtement le protocole dans le rapport}~: le
pendule bifilaire mesure l'inertie autour de l'axe \textbf{vertical} de la
configuration suspendue~; il faut donc suspendre la maquette de façon que
l'axe d'intérêt (tangage ou roulis) soit vertical, ou justifier
l'approximation faite.

\subsection{Le calage expérimental $K_{SYN}$}

La synthèse n'utilise pas directement $J_X$ et $J_Y$, mais~:
\begin{equation}
    J_{syn} = \frac{\sqrt{J_X\,J_Y}}{K_{SYN}} = \frac{0{,}848}{1{,}15}
            = 0{,}737\ \text{kg}\cdot\text{m}^2
    \quad\text{(les deux modes angulaires)},
    \qquad
    m_{syn} = \frac{m_{tot}}{1{,}15} = 15{,}70\ \text{kg}.
\end{equation}
Pourquoi~? L'identification en vol (méthode de la résonance,
§\ref{subsec:identres}) montre que les inerties \emph{effectives} vues par la
boucle diffèrent des valeurs pendulaires (dynamiques non modélisées~:
couplages, élasticité du rail, constante de force). La configuration
$K_{SYN} = 1{,}15$ correspond au meilleur vol réalisé~; elle est conservée
comme base expérimentale validée, et les ajustements ultérieurs se font
\emph{uniquement} par le choix des pôles (traçabilité~: un seul bouton de
réglage à la fois). C'est une démarche d'ingénieur parfaitement défendable, à
condition de la présenter comme telle~: modèle nominal $\to$ identification
$\to$ recalage $\to$ validation.

% ============================================================================
\section{Synthèse MIMO complète et détails d'implantation}
\label{sec:synthmimo}
% ============================================================================

\subsection{Chaîne de calcul par mode}

Pour chaque mode $j \in \{Z, T, R\}$, dans l'ordre~:
\begin{enumerate}
    \item modèle continu \eqref{eq:AB} avec l'inertie $J_j$ et
          $\tau_{act} = 5{,}565$~ms~;
    \item discrétisation exacte à $h = 40$~\textmu s (§\ref{sec:disc})~;
    \item placement des pôles $\{-80,-85,-90,-20\}$~rad/s $\to$ gains
          \texttt{LQI\_Q[j]}, \texttt{LQI\_QD[j]}, \texttt{LQI\_EPS[j]}
          (tableau~\ref{tab:gains})~;
    \item observateur, pôles $\{-700,-760,-820\}$~rad/s $\to$
          \texttt{OBS\_*[j]} (§\ref{sec:obs})~;
    \item vérification~: marge $|L(j\omega^*)| = 12{,}7 > 8$
          (§\ref{sec:marge}) et simulation de la boucle implantée
          (\texttt{simulate\_mode}).
\end{enumerate}

\subsection{Lissage de consigne}

La référence modale rejoint sa cible par un premier ordre de constante
$\tau_{ref} = 0{,}3$~s, réalisé en Euler~:
\begin{equation}
    \qm_{ref} \mathrel{+}= (\qm_{cible} - \qm_{ref})\cdot
    \underbrace{\frac{h}{\tau_{ref}}}_{\texttt{REF\_SMOOTH} = 1{,}333\e{-4}} .
\end{equation}
Rôle~: ne jamais demander d'échelon à une boucle dont les actionneurs
saturent vite~; c'est l'équivalent, côté consigne, du gain $K_w$ de mise en
forme du régime transitoire de la structure classique.

\subsection{Anti-windup en MIMO}
\label{subsec:antiwindupmimo}

Les saturations ont lieu sur les \emph{forces} ($0 \le F_k \le F_{max,k}$),
mais les intégrateurs vivent dans l'espace \emph{modal}. Le firmware reprojette
donc les forces saturées vers les modes~:
\begin{equation}
    \bar u = E\,F_{sat} - U_{stat},
\end{equation}
et compare mode par mode $\bar u_j$ à $u_{cmd,j}$~: l'intégrateur $j$ n'est mis
à jour que si $|u_{cmd,j} - \bar u_j| < \texttt{AW\_TOL}[j]$. Les tolérances
valent 2\,\% de l'effort statique, converties en couple par le bras de levier
moyen pour les modes angulaires~:
\begin{equation}
    \texttt{AW\_TOL} = \big[\,3{,}08\ \text{N},\;\; 0{,}431\ \text{N·m},\;\;
    0{,}400\ \text{N·m}\,\big].
\end{equation}
Le même $\bar u$ sert d'entrée à l'observateur \eqref{eq:obs}~: on informe le
modèle interne de ce qui a \emph{réellement} été appliqué, pas de ce qui a été
demandé.

% ============================================================================
\section{Identifications expérimentales}
\label{sec:ident}
% ============================================================================

\subsection{Inertie effective par résonance}
\label{subsec:identres}

Principe~: quand la boucle sature, elle oscille à une fréquence proche de
$\omega^{*}$ dont la valeur dépend de l'inertie réelle (§\ref{sec:marge}). En
mesurant la fréquence d'oscillation en vol avec le firmware SISO (gains
connus), on peut inverser numériquement la relation fréquence
$\leftrightarrow$ inertie~: le script simule la boucle pour une inertie
candidate, compare la fréquence dominante obtenue à la fréquence mesurée, et
itère par dichotomie (\texttt{inertia\_from\_resonance}). L'inertie de l'axe
s'en déduit via les bras de levier~: $J = m_{eff}\sum y_k^2$ (roulis) ou
$m_{eff}\sum x_k^2$ (tangage).

\subsection{Position du centre de gravité par les forces statiques}
\label{subsec:identcg}

En lévitation stationnaire, les forces moyennes $\bar F_k$ équilibrent le
poids \emph{et} les couples~: le centre de gravité est donc le barycentre des
forces~:
\begin{equation}
    x_{cg} = \frac{\sum_k x_k \bar F_k}{\sum_k \bar F_k},
    \qquad
    y_{cg} = \frac{\sum_k y_k \bar F_k}{\sum_k \bar F_k}.
\end{equation}
Avec les forces relevées en vol ($32{,}6$ / $44{,}8$ / $46{,}5$ / $31{,}8$~N),
les offsets obtenus sont inférieurs à 3~mm~: l'hypothèse \og couples statiques
nuls \fg{} de \texttt{U\_STAT} est justifiée. (Remarque~: la dissymétrie des
forces mesurées traduit surtout des différences de constante de force entre
inducteurs, que les intégrateurs compensent.)

% ============================================================================
\section{Récapitulatif~: d'où sort chaque constante}
% ============================================================================

\begin{table}[h]
    \centering
    \renewcommand{\arraystretch}{1.25}
    \begin{tabular}{lll}
        \toprule
        Constante & Formule & Référence \\
        \midrule
        \texttt{T\_MAT}    & $(G^TG)^{-1}G^T$ (moindres carrés)        & §\ref{subsec:T} \\
        \texttt{E\_MAT}    & $G^T$ (résultante et moments)             & §\ref{subsec:statique} \\
        \texttt{W\_MAT}    & $G(G^TG)^{-1} = T^T$ (norme minimale)     & §\ref{subsec:statique} \\
        \texttt{U\_STAT}   & $[m_{tot}g/K_{FF},\,0,\,0]$               & §\ref{subsec:statique} \\
        \texttt{F\_STAT}   & $W\,U_{stat} = 38{,}5$~N par inducteur    & §\ref{subsec:statique} \\
        \texttt{LQI\_Q/QD/EPS} & placement de pôles sur $(A_d, B_d)$, éq.~\eqref{eq:place} & §\ref{sec:placement} \\
        \texttt{LQI\_EPS\_INV} & $1/\texttt{LQI\_EPS}$ (init.\ sans à-coup de $\varepsilon$) & --- \\
        \texttt{OBS\_AD12} & $h$                                        & §\ref{sec:obs} \\
        \texttt{OBS\_AD23} & $-\tfrac{\tau_{act}}{J}(1-e^{-h/\tau_{act}}) \approx -h/J$ & §\ref{sec:obs} \\
        \texttt{OBS\_AD33} & $e^{-h/\tau_{act}} = 0{,}99284$            & §\ref{sec:disc} \\
        \texttt{OBS\_BD3}  & $1-e^{-h/\tau_{act}} = 0{,}00716$          & §\ref{sec:disc} \\
        \texttt{OBS\_L1/L2/L3} & placement dual sur $(A_{d,obs}^T, C^T)$ & §\ref{sec:obs} \\
        \texttt{AW\_TOL}   & $2\%$ de l'effort statique ($\times$ bras de levier) & §\ref{subsec:antiwindupmimo} \\
        \texttt{REF\_SMOOTH} & $h/\tau_{ref} = 1{,}333\e{-4}$           & §\ref{sec:synthmimo} \\
        \bottomrule
    \end{tabular}
    \caption{Traçabilité des constantes du firmware.}
\end{table}

\paragraph{Checklist \og refaire tout à la main \fg.}
\begin{enumerate}
    \item Mesurer la géométrie ($x_k$, $y_k$) et écrire $G$~;
          calculer $G^TG$ (diagonale par symétrie), puis $T$, $E = G^T$,
          $W = T^T$.
    \item Calculer $F_{stat} = m_{tot}g/(4K_{FF})$.
    \item Mesurer $J_X$, $J_Y$ au pendule bifilaire
          ($J = mgB^2T_p^2/4\pi^2L$)~; appliquer le calage $K_{SYN}$.
    \item Sommer les petites constantes de temps~: $\tau_{act}$
          (tableau~\ref{tab:tauact}).
    \item Écrire $A$, $B$ \eqref{eq:AB} pour chaque inertie~; discrétiser
          \eqref{eq:discret}.
    \item Placer les pôles (identification du polynôme caractéristique,
          comme dans le chapitre du rapport)~: gains $Q$, $Q_D$, $E_{PS}$.
    \item Observateur par dualité, pôles $\approx 10\times$ plus rapides.
    \item Vérifier $\omega^* = \sqrt{E_{PS}/Q_D}$ et la marge
          $Q\,Q_D/(J\,E_{PS}) > 8$.
    \item Simuler la boucle \emph{telle qu'implantée} (3 gains, observateur,
          coupe-bande, saturations).
    \item Essai en vol~; si oscillation croissante en saturation~: ralentir
          les pôles et reboucler en 6.
\end{enumerate}

\paragraph{Questions d'expert probables.}
\begin{itemize}
    \item \emph{Pourquoi 3 modes pour 4 actionneurs~?} --- Corps rigide, 3
          degrés de liberté verticaux~; la redondance (torsion) est dans le
          noyau de $E$ et n'est jamais excitée par $W$.
    \item \emph{LQI = optimal~?} --- Structure LQI, gains par placement de
          pôles~; pas de critère quadratique résolu ici.
    \item \emph{Pourquoi l'intégrateur~?} --- Erreur statique nulle malgré
          $K_{FF}$, constantes de force incertaines et masse variable.
    \item \emph{Pourquoi la marge doit-elle être $>8$ et pas $>1$~?} ---
          Boucle conditionnellement stable~: la saturation \emph{réduit} le
          gain effectif~; 12{,}7 laisse $\approx 22$~dB de réduction admissible.
    \item \emph{D'où viennent les inerties~?} --- Pendule bifilaire + calage
          par identification en vol ($K_{SYN}$).
\end{itemize}

\end{document}
